% Speaker notes: one frame per way a note is written (docs/speaker-notes.md).
% !variants = notes notes-right notes-plain notes-compressed notes-xelatex
% The plain build has no note pages at all (what a person compiles every day); the variants show
% them as beamer's options and templates do (tests/decks/build.py VARIANTS). What each note must
% read as is tests/test_speaker_notes.py's EXPECTED.
\documentclass{beamer}
\title{Speaker Notes}
\author{Tester}
\date{}

\begin{document}

\begin{frame}{A plain sentence}
  The note under this frame is one sentence.
  \note{Remember to greet the audience before the first slide.}
\end{frame}

\begin{frame}{Several paragraphs}
  Three paragraphs of notes.
  \note{The first paragraph sets the scene for the talk and says why the topic matters to the
    people in the room today.

    The second paragraph gives the one number to remember.

    The third paragraph is a short reminder to slow down.}
\end{frame}

\begin{frame}{Lists in a note}
  A bullet list and a numbered list.
  \note{Points to make:
    \begin{itemize}
      \item the parser reads the PDF
      \item the classifier groups the lines
    \end{itemize}
    Then in order:
    \begin{enumerate}
      \item compile
      \item convert
      \item sync
    \end{enumerate}}
\end{frame}

\begin{frame}{Note items}
  Three notes given as items.
  \note[item]{Ask who has used beamer before.}
  \note[item]{Show the converted deck.}
  \note[item]{Mention the speaker notes.}
\end{frame}

\begin{frame}{Bold, italic and emphasis}
  Styles inside a note.
  \note{This is \textbf{bold}, this is \textit{italic} and this is \emph{emphasised}; a word
    \textbf{half}way styled stays one word.}
\end{frame}

\begin{frame}{Inline math}
  Formulas inside a note.
  \note{The circle is $x^2 + y^2 = r^2$, the angle is $\alpha$ and $a_i \le b_i$ for every $i$.}
\end{frame}

\begin{frame}{Accents and quotes}
  Letters beyond ASCII.
  \note{Café, naïve, Müller, Straße and señor --- “curly quotes” and ``TeX quotes'', 1--2 pages,
    an efficient office workflow.}
\end{frame}

\begin{frame}{Links}
  A link and a URL.
  \note{See \href{https://example.com/talk}{the talk page} and \url{https://example.org/notes}.}
\end{frame}

\begin{frame}{A long note}
  The note under this frame is longer than one note page.
  \note{This note is long on purpose. It runs past the bottom of its note page, because beamer
    sets every note into one page and never breaks it onto a second one. The words that run off
    the page are still in the PDF, drawn below its edge, and a person who wrote them wants to see
    them in the speaker notes all the same. So the reader takes every word the note holds, on the
    page or below it. The rest of this note is filler that says the same thing again and again
    until it is long enough to overflow. One two three four five six seven eight nine ten eleven
    twelve thirteen fourteen fifteen sixteen seventeen eighteen nineteen twenty. One two three
    four five six seven eight nine ten eleven twelve thirteen fourteen fifteen sixteen seventeen
    eighteen nineteen twenty. One two three four five six seven eight nine ten eleven twelve
    thirteen fourteen fifteen sixteen seventeen eighteen nineteen twenty. One two three four five
    six seven eight nine ten eleven twelve thirteen fourteen fifteen sixteen seventeen eighteen
    nineteen twenty. One two three four five six seven eight nine ten eleven twelve thirteen
    fourteen fifteen sixteen seventeen eighteen nineteen twenty. One two three four five six
    seven eight nine ten eleven twelve thirteen fourteen fifteen sixteen seventeen eighteen
    nineteen twenty. One two three four five six seven eight nine ten eleven twelve thirteen
    fourteen fifteen sixteen seventeen eighteen nineteen twenty. One two three four five six
    seven eight nine ten eleven twelve thirteen fourteen fifteen sixteen seventeen eighteen
    nineteen twenty. One two three four five six seven eight nine ten eleven twelve thirteen
    fourteen fifteen sixteen seventeen eighteen nineteen twenty. One two three four five six
    seven eight nine ten eleven twelve thirteen fourteen fifteen sixteen seventeen eighteen
    nineteen twenty. The last sentence of the long note.}
\end{frame}

\begin{frame}{Overlays}
  \begin{itemize}
    \item First step
    \pause
    \item Second step
    \pause
    \item Third step
  \end{itemize}
  \note{A note on every step.}
  \note<2>{Only on the second step.}
\end{frame}

\begin{frame}{No note}
  This frame has no note.
\end{frame}

\begin{frame}{A note after the frame}
  Its note is written after the frame ends.
\end{frame}
\note{Written outside the frame, after it.}

\begin{frame}[noframenumbering]{Not numbered}
  This frame does not count.
  \note{A note on a frame that takes no number.}
\end{frame}

\begin{frame}{The last frame}
  The end.
  \note{Thank the audience.}
\end{frame}

\end{document}
